\section{Fixed-degree row formulation}

Let $K=2m+1\ge7$ be an odd prime, let
$\zeta=\exp(2\pi i/K)$, and put $\tau=\zeta+\zeta^{-1}$. The full subgroup
trace polynomial is $H_K(T)=(T-2)C_K(T)$, where $C_K$ is the minimal
polynomial of $\tau$. For a normalized row ratio $r\ne1$ the affine trace map is
\[
Y=\alpha_rX+\beta_r,\qquad
\alpha_r=U_{r-1}(\tau/2),\qquad\beta_r=-2U_{r-2}(\tau/2).
\]
Removing the shared invalid characteristic-zero row trace gives fixed-degree
cofactors $P_{K,r},Q_{K,r}\in\mathbf Z[\tau][T]$ of degree $m$.
At every split prime $p\equiv1\pmod K$ and embedding prime $\mathfrak q\mid p$,
the leading coefficients remain nonzero and
\begin{equation}\label{frontier:eq:gcd-count}
I_{1,r}(\mathfrak q)=\deg\gcd(P_{K,r},Q_{K,r}\bmod\mathfrak q).
\end{equation}
This fixed-degree statement is the specialization safeguard. Older
projective dehomogenizations may lose degree at infinity and remain only
one-way obstructions.

\section{Principal subresultants are sufficient}
For fixed-degree polynomials $P,Q$ let
$\operatorname{psc}_j(P,Q)=[T^j]S_j(P,Q;T)$.
The standard subresultant theorem \cite{BT1971} states that, for
$g=\deg\gcd(P,Q)$, the first nonzero subresultant is a nonzero scalar
multiple of the gcd, of degree $g$.
\begin{lemma}[Principal coefficients detect gcd degree]\label{frontier:lem:psc}
For $1\le d\le\min(\deg P,\deg Q)$,
\[
\deg\gcd(P,Q)\ge d\iff
\operatorname{psc}_0=\cdots=\operatorname{psc}_{d-1}=0.
\]
\end{lemma}
\begin{proof}
If the gcd degree is at least $d$, these subresultants vanish. Otherwise its
degree $g<d$ gives a first nonzero subresultant of degree exactly $g$, so its
principal coefficient is nonzero.
\end{proof}
Set $\pi_j=\operatorname{psc}_j(P_{K,r},Q_{K,r})$ for $j=0,1,2$, and define
\[
J^{\rm psc}_{K,r}=(\pi_0,\pi_1,\pi_2),\qquad
J^{\rm psc}_{K,r}\cap\mathbf Z=\Pi_{K,r}^{(3)}\mathbf Z.
\]
\begin{theorem}[Three-scalar exact split support]\label{frontier:thm:psc-support}
For every split prime $p\equiv1\pmod K$,
\[
p\mid\Pi_{K,r}^{(3)}\iff
\text{some }\mathfrak q\mid p\text{ has }I_{1,r}(\mathfrak q)\ge3.
\]
\end{theorem}
\begin{proof}
Use \eqref{frontier:eq:gcd-count} and Lemma~\ref{frontier:lem:psc}.
A rational prime divides the contraction precisely when the ideal is contained
in some prime ideal above it.
\end{proof}
This is equality of split prime support, not of integral valuations.
Exact finite Smith-normal-form controls are
\[
\begin{array}{c|c|c}
(K,r)&\Pi_{K,r}^{(3)}&\text{prime-to-}2K\text{ part}\\\hline
(7,2)&7&1\\
(11,7)&11^2&1\\
(13,3)&13^3&1\\
(13,12)&2^4\,5^3\,13^3&5^3\\
(17,10)&2^2\,17^5\,103&103.
\end{array}
\]
At $(13,12)$ the full/shift support instead has prime-to-$2K$ valuation $5^2$.

\section{The intrinsic Sylvester/Fitting locus}
Let
\[
\mathcal S_{P,Q}:k[T]_{<m}\oplus k[T]_{<m}\longrightarrow k[T]_{<2m},
\qquad (U,V)\longmapsto UP+VQ.
\]
\begin{theorem}[Sylvester nullity]\label{frontier:thm:nullity}
For nonzero degree-$m$ polynomials,
$\dim\ker\mathcal S_{P,Q}=\deg\gcd(P,Q)$.
\end{theorem}
\begin{proof}
Write $P=GP_1,Q=GQ_1$, with $(P_1,Q_1)=1$ and $g=\deg G$.
Kernel vectors are $(WQ_1,-WP_1)$ with $\deg W<g$.
\end{proof}
Thus gcd degree at least three is equivalent to rank at most $2m-3$, or to
vanishing of all $(2m-2)$-minors. At gcd degree two the corresponding compound
matrix has rank one. These are intrinsic determinantal descriptions of the
same field-valued support.

\section{Shift-resultant transversality}
The squarefreeness hypothesis must be stated on the correct input. A
squarefree gcd alone does not suffice: $f=T^2$, $g=T(T-1)$ have gcd degree
one, but $\Res(f,g+Z)=Z^2$.

\begin{theorem}[Universal initial resultant coefficient]\label{frontier:thm:shift-psc}
Let $f=GA,g=GB$ have degrees $n,m$, with $G$ monic of degree $d$.
For degree-preserving perturbations $f_Z=f+Z\phi$, $g_Z=g+Z\psi$, put
$D=A\psi-B\phi$. With the root-product resultant convention,
\[
\Res(f_Z,g_Z)=(-1)^{d(n-d)}\Res(A,B)\Res(G,D)Z^d+O(Z^{d+1}).
\]
This coefficient identity does not require $G$ to be squarefree. When
$G=\gcd(f,g)$ and $\Res(G,D)\ne0$, the order is exactly $d$.
\end{theorem}
\begin{proof}
For generic characteristic-zero coefficients, assume $G$ squarefree and
coprime to $A$. A root $y$ of $G$ contributes $ZD(y)/A(y)+O(Z^2)$ to the
resultant root product. All other roots contribute
$\Res(A,G)\Res(A,B)$, and
$\Res(A,G)=(-1)^{d(n-d)}\prod_{G(y)=0}A(y)$.
Multiplication proves the displayed identity and vanishing of every lower
coefficient. These are polynomial identities over the integers in the
coefficients of monic $G$ and of $A,B,\phi,\psi$. Therefore specialization
proves them in arbitrary characteristic and on multiple-root strata.
\end{proof}

In particular, for monic same-degree $f,g$, squarefree $G=\gcd(f,g)$,
$\gcd(A,G)=1$, and $\psi$ nonzero at every root of $G$, a fixed-$f$ shift gives
\[
\ord_{Z=0}\Res(f,g+Z\psi)=d,\qquad
[Z^d]\Res(f,g+Z\psi)=\Res(A,B)\Res(A,G)\Res(G,\psi).
\]
Up to the fixed subresultant sign, $\Res(A,B)$ is the first nonzero PSC.
The condition $\gcd(A,G)=1$ says that the common roots are simple in $f$;
squarefreeness of $f$ is a stronger sufficient condition. Without it the
coefficient formula still holds but can vanish.

For the row shift put
$q_Z(T)=\alpha_r^{-m}R_r(\alpha_rT+\beta_r+Z)$, so
$\psi(T)=\alpha_r^{-m}R_r'(\alpha_rT+\beta_r)$.
On the maintained split locus both row cofactors are squarefree and
$\alpha_r\ne0$. Hence $q_0'(y)=\alpha_r\psi(y)\ne0$ at common roots.
Writing $P=GA,q_0=GB$ gives
\begin{equation}\label{frontier:eq:explicit-local}
[Z^d]\Res(P,q_Z)
 =\pm\alpha_r^{-d}\operatorname{psc}_d(P,q_0)
 \Disc(G)\Res(A,G)\Res(G,B).
\end{equation}
Every multiplier is nonzero. The first three shift coefficients therefore
vanish precisely on the order-three PSC locus, with no conclusion about
equality of integral valuations.

\section{Inverse resultant transversality}
\begin{theorem}[Simple-gcd derivative]\label{frontier:thm:derivative}
Let $\gcd(f,g)=Y-y$, write $f=(Y-y)A,g=(Y-y)B$, and preserve the displayed
degrees in $f_Z=f+Z\phi,g_Z=g+Z\psi$. Then
\[
\left.\frac{d}{dZ}\Res(f_Z,g_Z)\right|_0
=(-1)^{\deg f-1}\Res(A,B)
       \bigl(f'(y)\psi(y)-g'(y)\phi(y)\bigr).
\]
A root may be multiple in one input; simplicity in both is not required.
\end{theorem}
\begin{proof}
Apply Theorem~\ref{frontier:thm:shift-psc} with $G=Y-y$.
\end{proof}

For $F(X)=E(X^2)+XO(X^2)$, translation by $Z/2$ has first variations
\[
\phi(Y)=\tfrac12(O(Y)+2YO'(Y)),\qquad\psi(Y)=E'(Y).
\]
Suppose the characteristic is not two, $F$ is nonconstant and squarefree,
and $F(0)\ne0$. Let $G=\gcd(E,O)$ be monic, with $E=GA,O=GB$.
Then $G$ is squarefree and its roots $y$ are nonzero. At a common root,
\[
D(y)=G'(y)(A(y)^2-yB(y)^2),
\]
and for $z^2=y$,
\[
F'(z)F'(-z)=-4yG'(y)^2(A(y)^2-yB(y)^2)\ne0.
\]
Thus $D$ is invertible modulo $G$.

To include all degree-drop cases, use the fixed ambient parity degrees in the
Sylvester matrix. At least one of $E,O$ has its full ambient degree, because
$F$ has its displayed degree. Swap the inputs if needed. For a full-degree
first input $f$ and a second input $g$ of degree at most its ambient degree,
the Sylvester kernel is still $(BW,-AW)$, $\deg W<d$, and its cokernel is
identified with $k[Y]/(G)$. The derivative on kernel to cokernel is
multiplication by $B\phi-A\psi=-D$ (up to the input swap).
Its determinant is nonzero. Constant invertible row and column operations
reduce the undeformed Sylvester matrix to an identity block and a zero
$d$-block; the Schur complement then shows its determinant has order exactly
$d$. This also covers $O=0$, where the same kernel description uses $B=0$.
Consequently, with the fixed-degree convention,
\[
\ord_{Z=0}\rho(Z)=\deg\gcd(E,O),\qquad \rho(Z)=\Res(E_Z,O_Z),
\]
and
\[
\deg\gcd(E,O)\ge2\iff\rho(0)=\rho'(0)=0.
\]
For the unique-common-root formula when the displayed input degrees are
preserved, the local factor is
\[
E'(y)^2-yO'(y)^2=-\frac{F'(z)F'(-z)}{4y}.
\]
The example $E=(Y-1)^2,O=Y-1$ over $\mathbf Q$ has $E'(1)=0$ and squarefree
$F=(X^2-1)(X^2+X-1)$ with $F(0)=1$; its translated resultant has linear
coefficient one. It is a finite control for the repaired multiple-input-root
argument. These local identities do not bound the largest split prime in the
integral support and do not establish microscopic BGS universality.
